\scanpage{10}\mspage{\handwritten{17}}
$H^{3}(Y,\underline G_{\mathrm m})\handwritten{\simeq}\underline{\mathrm Q}/\underline{\mathrm Z}$, d'où le corollaire.

Soit \deleted{\unsure{$i$:}}
\[
  i:\handwritten{\unsure{\eta|y}\longrightarrow}Y
\]
l'inclusion du point générique, et posons
\[
  \underline B=i_{\handwritten{\unsure{\times|\ast}}}(i^{*}(P)) \quad , \tag*{(4.8)}
\]
de sorte qu'on a un homomorphisme canonique
\[
  P\handwritten{\longrightarrow}\underline B \quad , \tag*{(4.9)}
\]
et on posera
\[
  \underline E=\operatorname{Ker}(P\handwritten{\to}\underline B)
  \quad ,\quad \underline F=\operatorname{Coker}(P\handwritten{\to}\underline B)\quad . \tag*{(4.10)}
\]

Les faisceaux $\underline E$ et $\underline F$ sont des ``faisceaux gratte-ciel'' , qui sont donc connus lorsqu'on connait leurs restrictions aux points fermés de $Y$. Leur\handwritten{s} fibre\handwritten{s} géométriques au dessus d'un tel point $y$ s'obtiennent comme les noyaux et conoyau de l'homomorphisme
\[
  \deleted{\unsure{\mathrm{Pic}(}}\inserted{\mathrm{Pic}(\overset{\handwritten{\sim}}{X})=}
  \mathrm{Pic}(\overset{\handwritten{\sim}}{X}/\overset{\handwritten{\sim}}{Y})\handwritten{\longrightarrow}
  \mathrm{Pic}(\overset{\handwritten{\sim}}{X}_{\handwritten{\tilde\eta}}/\handwritten{\tilde\eta})\quad ,
\]
où les $\handwritten{\sim}$ dénotent l'effet de la localisation stricte relativement à un point géométrique $\bar y$ au dessus de $y$. \deleted{\unsure{Lorsque}} Compte tenu que \deleted{L}\handwritten{l}e morphisme précédent se factorise par le sous-groupe $\mathrm{Pic}(\overset{\handwritten{\sim}}{X}_{\handwritten{\eta}})$, et que $\mathrm{Pic}(\overset{\handwritten{\sim}}{X})\handwritten{\longrightarrow}\mathrm{Pic}(\overset{\handwritten{\sim}}{X}_{\handwritten{\tilde\eta}})$ est surjectif puisque $\overset{\handwritten{\sim}}{X}$ est régulier, on trouve
\[
  \deleted{\underline{\unsure{E|F}}_{\bar y}}\qquad
  \underline F_{\bar{\deleted{\unsure{y}}\handwritten{y}}\deleted{\unsure{*}}}=
  \mathrm{Pic}(\overset{\handwritten{\sim}}{X}_{\handwritten{\eta}}/\handwritten{\tilde\eta})\,/\,\mathrm{Pic}(\overset{\handwritten{\sim}}{X}_{\handwritten{\tilde\eta}}) \tag*{(4.11)}
\]
\[
  \underline E_{\bar y}\deleted{\unsure{\delta_y|y}}=
  \operatorname{Ker}\bigl(\mathrm{Pic}(\overset{\handwritten{\sim}}{X})\handwritten{\longrightarrow}\mathrm{Pic}(\overset{\handwritten{\sim}}{X}_{\handwritten{\tilde\eta}})\bigr)\quad . \tag*{(4.12)}
\]
Compte tenu que $\mathrm{Br}(\handwritten{\tilde\eta})=0$ si $k(y)$ est parfait, et que c'est un groupe de $p(y)$-torsion en tous cas, où $p(y)$ est l'exposant cara\deleted{\unsure{tc}}\handwritten{ct}éristique de $k(y)$, on trouve que $\underline F_{\bar y}$ est un groupe de $p(y)$-torsion, nul si $k(y)$ est parfait, Désignant par $\handwritten{\delta}_{y},\handwritten{\delta'}_{y},\handwritten{\delta''}_{y}$ les entiers analogues à $\handwritten{\delta},\handwritten{\delta'},\handwritten{\delta''}$ pour la situation localisée $\overset{\handwritten{\sim}}{X}\handwritten{\longrightarrow}\overset{\handwritten{\sim}}{Y}$, un argument déjà \deleted{\unsure{sig}} fait plus haut nous montre que $\underline F_{\bar y}$ \deleted{est nul \unsure{également}} est nul également si $\handwritten{\delta}_{y}=1$, par exemple si $\overset{\handwritten{\sim}}{X}$ admet un\unsure{e} section sur $\overset{\handwritten{\sim}}{Y}$. Nous allons supposer pour simplifier, par la suite, que
% archive sheet number: 9
\clearpage
